\documentclass[11pt]{article}
\usepackage[T1]{fontenc}
\usepackage[utf8]{inputenc}
\usepackage{amsmath,amssymb,amsthm,mathtools}
\usepackage{stmaryrd}
\usepackage[margin=1.1in]{geometry}

\theoremstyle{definition}
\newtheorem{definition}{Definition}[section]
\theoremstyle{plain}
\newtheorem{proposition}[definition]{Proposition}
\newtheorem{lemma}[definition]{Lemma}
\newtheorem{theorem}[definition]{Theorem}
\theoremstyle{remark}
\newtheorem{remark}[definition]{Remark}

\newcommand{\N}{\mathbb{N}}
\newcommand{\Terms}{\mathcal{T}}
\newcommand{\Msg}{\mathcal{M}}
\newcommand{\Msgf}{\Msg_{\bot}}
\newcommand{\Fr}{\mathrm{Fr}}
\newcommand{\eqE}{\equiv_{E}}
\newcommand{\View}{\mathcal{V}}
\newcommand{\view}{\operatorname{view}}
\newcommand{\Ops}{\mathrm{Ops}}
\newcommand{\Pub}{\mathrm{Pub}}
\newcommand{\Dt}{\mathrm{Dt}}
\newcommand{\Act}{\mathrm{Act}}
\newcommand{\Proc}{\mathrm{Proc}}
\newcommand{\aout}{\mathsf{out}}
\newcommand{\ain}{\mathsf{in}}
\newcommand{\anew}{\mathsf{new}}
\newcommand{\aev}{\mathsf{ev}}
\newcommand{\acomm}{\mathsf{comm}}
\newcommand{\known}{\operatorname{known}}
\newcommand{\knows}{\pv{knows}}
\newcommand{\Fact}{\mathrm{Fact}}
\newcommand{\Closed}{\mathrm{Closed}}
\newcommand{\nonces}{\operatorname{nonces}}
\newcommand{\comms}{\operatorname{comms}}
\newcommand{\mrg}{\operatorname{merge}}
\newcommand{\Enemy}{\mathrm{Enemy}}
\newcommand{\Attacks}{\operatorname{Attacks}}
\newcommand{\Secret}{\operatorname{Secret}}
\newcommand{\Corr}{\operatorname{Corr}}
\newcommand{\sem}[1]{\llbracket #1 \rrbracket}
\newcommand{\pv}[1]{\mathsf{#1}}
\newcommand{\fail}{\pv{fail}}
\newcommand{\true}{\pv{true}}
\newcommand{\ev}{\Downarrow}
\newcommand{\Npub}{\mathcal{N}_{\mathrm{pub}}}
\newcommand{\Npriv}{\mathcal{N}_{\mathrm{priv}}}
\newcommand{\fn}{\operatorname{fn}}
\newcommand{\obs}{\operatorname{obs}}
\newcommand{\Tr}{\operatorname{Tr}}
\newcommand{\attacker}{\pv{attacker}}
\newcommand{\event}{\pv{event}}
\newcommand{\interleaves}{\operatorname{Interleaves}}
\newcommand{\red}[1]{\text{(Red #1)}}
\newcommand{\CQ}{\mathcal{C}_Q}
\newcommand{\Azero}{\mathcal{A}_0}

\title{A trace semantics for the process calculus of ProVerif}
\author{}
\date{}

\begin{document}
\maketitle

\section{The trace semantics}

\subsection{Messages}

\begin{definition}[Signature, terms]
A \emph{signature} is a family $\Sigma=(\Sigma_n)_{n\in\N}$ of sets; $f\in\Sigma_n$ is a constructor of arity $n$. The set $\Terms=\Terms(\Sigma)$ of \emph{raw terms} is generated by
\[
t \Coloneqq \nu_k \;(k\in\N) \;\mid\; f(t_1,\dots,t_n)\;(f\in\Sigma_n).
\]
A \emph{substitution} is a map $\sigma:\N\to\Terms$; $t\sigma$ replaces each $\nu_k$ in $t$ by $\sigma(k)$.
\end{definition}

\begin{definition}[Theory]
A \emph{theory} over $\Sigma$ is a pair $(E,\Fr)$ with $E\subseteq\Terms\times\Terms$ and $\Fr\subseteq\bigcup_n\Sigma_n$ such that
\begin{enumerate}
\item[(H)] for every $(l,r)\in E$, both $l$ and $r$ are of the form $f(\dots)$ with $f\notin\Fr$.
\end{enumerate}
The relation $\eqE$ is the least congruence on $\Terms$ containing $l\sigma\eqE r\sigma$ for all $(l,r)\in E$ and all substitutions $\sigma$. The set of \emph{messages} is $\Msg=\Terms/{\eqE}$, with $[t]$ the class of $t$. Constructors act on $\Msg$ by $f([t_1],\dots,[t_n])=[f(t_1,\dots,t_n)]$, and $\nu_k$ also denotes $[\nu_k]$.
\end{definition}

A $\Sigma$-algebra $A$ \emph{satisfies} $E$ if $l\sigma_A=r\sigma_A$ for every $(l,r)\in E$ and every valuation $\sigma_A:\N\to A$, where $t\sigma_A$ is the evaluation of $t$ in $A$.

\begin{proposition}
$\Msg$ satisfies $E$. For every $\Sigma$-algebra $A$ satisfying $E$ and every $v:\N\to A$ there is a unique homomorphism $\Msg\to A$ sending $\nu_k$ to $v(k)$.
\end{proposition}

\begin{definition}[View]
Let $\View=\N\;\uplus\;\{f(m_1,\dots,m_n)\mid f\in\Fr\cap\Sigma_n,\ m_i\in\Msg\}\;\uplus\;\{\bot\}$. Define $\view_0:\Terms\to\View$ by
\[
\view_0(\nu_k)=k,\qquad
\view_0(f(t_1,\dots,t_n))=
\begin{cases}
f([t_1],\dots,[t_n]) & f\in\Fr,\\
\bot & f\notin\Fr.
\end{cases}
\]
\end{definition}

\begin{proposition}\label{prop:view}
The following hold.
\begin{enumerate}
\item $s\eqE t$ implies $\view_0(s)=\view_0(t)$; hence $\view_0$ induces $\view:\Msg\to\View$.
\item $\view(f(m_1,\dots,m_n))=f(m_1,\dots,m_n)$ if $f\in\Fr$, and $\bot$ otherwise.
\item If $\view(m)=f(m_1,\dots,m_n)$ then $m=f(m_1,\dots,m_n)$.
\item $k\mapsto\nu_k$ is injective $\N\to\Msg$.
\end{enumerate}
\end{proposition}

\begin{definition}[May-fail messages, lifting]
Let $\Msgf=\Msg\uplus\{\bot\}$. For $g:\Msg^n\to\Msgf$ let $g^{\uparrow}:\Msgf^n\to\Msgf$ be $g^{\uparrow}(\vec u)=g(\vec u)$ if every $u_i\in\Msg$ and $\bot$ otherwise; $f^{\uparrow}$ for $f\in\Sigma_n$ is the lift of $f:\Msg^n\to\Msg\subseteq\Msgf$. A \emph{destructor} of arity $n$ is a function $\Msgf^n\to\Msgf$.
\end{definition}

\begin{definition}[Destructors from rewrite rules]\label{def:rules}
Fix disjoint countable sets of \emph{ordinary variables} $x$ and \emph{may-fail variables} $u$. Patterns over $\Fr$ are generated by $p\Coloneqq x\mid u\mid f(p_1,\dots,p_n)$ with $f\in\Fr\cap\Sigma_n$. A \emph{matching} of $\vec u\in\Msgf^n$ against $\vec p$ is a map $\theta$ from the variables of $\vec p$ to $\Msgf$, sending ordinary variables into $\Msg$, such that $p_i\theta=u_i$ for all $i$, where $f(p_1,\dots,p_n)\theta=f(p_1\theta,\dots,p_n\theta)$ if every $p_j\theta\in\Msg$ and is undefined otherwise. By Proposition~\ref{prop:view}, at most one matching exists.

A \emph{right-hand side} is $q\Coloneqq\fail\mid u\mid M$, where $M$ is a term over constructors and ordinary variables; $q\theta\in\Msgf$ is $\bot$, $\theta(u)$, and the evaluation of $M$ in $\Msg$ under $\theta$ respectively. A list of rules $g(\vec p_1)\to q_1;\ \dots;\ g(\vec p_k)\to q_k$, each variable of $q_i$ occurring in $\vec p_i$, defines the destructor
\[
g(\vec u)=\begin{cases}
q_j\theta & j \text{ least such that a matching } \theta \text{ of } \vec u \text{ against } \vec p_j \text{ exists},\\
\bot & \text{if there is none.}
\end{cases}
\]
\end{definition}

\begin{definition}[Public operations, deduction]
Let $\Pub\subseteq\bigcup_n\Sigma_n$ and let $\Dt$ be a set of destructors. Set $\Ops=\{f^{\uparrow}\mid f\in\Pub\}\cup\Dt$. A set $S\subseteq\Msg$ is \emph{$\Ops$-closed} if for every $o\in\Ops$ of arity $n$ and every $\vec u\in(S\uplus\{\bot\})^n$, $o(\vec u)\in\Msg$ implies $o(\vec u)\in S$. For $K\subseteq\Msg$,
\[
\langle K\rangle=\bigcap\{S\supseteq K\mid S\text{ is $\Ops$-closed}\},
\qquad
K\vdash m\iff m\in\langle K\rangle .
\]
\end{definition}

\begin{proposition}
$\langle\cdot\rangle$ is a closure operator on $\mathcal{P}(\Msg)$, $\langle K\rangle$ is $\Ops$-closed, and $K\vdash m_1,\dots,K\vdash m_n$ imply $K\vdash f(m_1,\dots,m_n)$ for $f\in\Pub\cap\Sigma_n$. In particular $\emptyset\vdash a$ for $a\in\Pub\cap\Sigma_0$.
\end{proposition}

\subsection{Traces and processes}

\begin{definition}[Actions, traces]
\[
\Act=\{\aout(c,m),\ \ain(c,m),\ \acomm(c,m),\ \anew(k),\ \aev(m)\mid c,m\in\Msg,\ k\in\N\}.
\]
A \emph{trace} is a word $t\in\Act^*$; $\varepsilon$ is the empty trace and $l\le t$ means $l$ is a prefix of $t$. The actions $\aout(c,m)$ and $\ain(c,m)$ are \emph{unfinished}; $\Closed(t)$ holds if no action of $t$ is unfinished. $\nonces(t)\in\N^*$ is the image of $t$ under the monoid homomorphism with $\anew(k)\mapsto k$ and every other action $\mapsto\varepsilon$. A \emph{process} is a subset of $\Act^*$; $\Proc=\mathcal{P}(\Act^*)$.
\end{definition}

\begin{definition}[Sequential constructors]\label{def:seq}
For $a\in\Act$, $P\in\Proc$, $F:\Msg\to\Proc$:
\begin{align*}
a\cdot P &= \{\varepsilon\}\cup\{a\,t\mid t\in P\}, &
\mathbf{0} &= \{\varepsilon\},\\
\aout(c,m).P &= \aout(c,m)\cdot P, &
\aev(m).P &= \aev(m)\cdot P,\\
\ain(c).F &= \textstyle\bigcup_{m\in\Msg}\ain(c,m)\cdot F(m), &
\anew.F &= \textstyle\bigcup_{k\in\N}\anew(k)\cdot\{t\in F(\nu_k)\mid k\notin\nonces(t)\}.
\end{align*}
\end{definition}

\begin{definition}[Parallel composition]\label{def:par}
Let $I$ be a set. For $s\in(I\times\Act)^*$ and $i\in I$, $s{\restriction}i\in\Act^*$ is the word of actions of $s$ tagged $i$, in order, and $\overline{s}\in\Act^*$ is $s$ with tags removed. For $f:I\to\Act^*$, $\interleaves(f,s)$ holds if $s{\restriction}i=f(i)$ for all $i$. The relation $\comms\subseteq(I\times\Act)^*\times\Act^*$ is the least one such that
\begin{align*}
&\comms(\varepsilon,\varepsilon),\\
&\comms(s,t)\Rightarrow\comms\big((i,a)\,s,\ a\,t\big),\\
&\comms(s,t)\wedge i\ne j\Rightarrow\comms\big((i,\aout(c,m))\,(j,\ain(c,m))\,s,\ \acomm(c,m)\,t\big).
\end{align*}
For $P:I\to\Proc$,
\begin{multline*}
\mrg_I(P)=\{t\mid \exists f,s.\ (\forall i.\ f(i)\in P(i))\ \wedge\ \interleaves(f,s)\\
\wedge\ \comms(s,t)\ \wedge\ \nonces(t)\text{ has no repeated letter}\}.
\end{multline*}
$P\mid Q=\mrg_{\{0,1\}}(0\mapsto P,\ 1\mapsto Q)$ and $\,!P=\mrg_{\N}(k\mapsto P)$.
\end{definition}

\begin{proposition}\label{prop:comms}
Let $\comms(s,t)$.
\begin{enumerate}
\item $\nonces(t)=\nonces(\overline{s})$.
\item If $a$ occurs in $t$ then $(i,a)$ occurs in $s$ for some $i$, or $a=\acomm(c,m)$ and $(i,\aout(c,m))$, $(j,\ain(c,m))$ occur in $s$ for some $i\ne j$.
\item If $\Closed(t)$ and $(k,\ain(c,m))$ occurs in $s$, then $(i,\aout(c,m))$ occurs in $s$ for some $i\ne k$.
\item $!P=P\mid{!P}$ for every $P\in\Proc$.
\end{enumerate}
\end{proposition}

\subsection{Translation of processes}

\begin{definition}[Source syntax]\label{def:syntax}
Fix a set of variables $x,y,\dots$ and a set of names $a,b,c,\dots$ containing the free names $\Npub\uplus\Npriv$. Terms, expressions and processes are generated by
\begin{align*}
M,N &\Coloneqq x\mid a\mid f(M_1,\dots,M_n),\\
D &\Coloneqq M\mid h(D_1,\dots,D_n)\mid\fail,\\
P,Q &\Coloneqq \mathbf{0}\mid \pv{out}(N,M);P\mid \pv{in}(N,x);P\mid P\mid Q\mid\,!P\mid\pv{new}\ a;P\\
&\qquad\mid \pv{let}\ x=D\ \pv{in}\ P\ \pv{else}\ Q\mid \pv{if}\ M\ \pv{then}\ P\ \pv{else}\ Q\mid \pv{event}(M);P,
\end{align*}
where $f$ ranges over constructors and $h$ over constructors and destructors. Types are erased.
\end{definition}

\begin{definition}[Model]\label{def:model}
A \emph{model} consists of a signature $\Sigma$ whose nullary symbols include $\Npub\uplus\Npriv$, a theory $(E,\Fr)$ over $\Sigma$, a finite set of destructors each given by rules as in Definition~\ref{def:rules}, sets $\Pub$, $\Dt$ of public constructors and public destructors, and a closed process $P_0$ with $\fn(P_0)\subseteq\Npub\uplus\Npriv$.
\end{definition}

\begin{definition}[Translation]
An \emph{environment} $\rho$ maps variables and names to $\Msg$; $\rho_0$ maps each free name $a$ to the constant $a\in\Msg$. Write $\rho M\in\Msg$ for the homomorphic extension. Expressions are evaluated in $\Msgf$:
\[
\sem{M}\rho=\rho M,\qquad \sem{\fail}\rho=\bot,\qquad
\sem{f(\vec D)}\rho=f^{\uparrow}(\sem{\vec D}\rho),\qquad
\sem{g(\vec D)}\rho=g(\sem{\vec D}\rho)
\]
for a constructor $f$ and a destructor $g$. Processes are translated by
\begin{align*}
\sem{\mathbf{0}}\rho&=\mathbf{0},&
\sem{\pv{out}(N,M);P}\rho&=\aout(\rho N,\rho M).\sem{P}\rho,\\
\sem{\pv{event}(M);P}\rho&=\aev(\rho M).\sem{P}\rho,&
\sem{\pv{in}(N,x);P}\rho&=\ain(\rho N).\big(m\mapsto\sem{P}\rho[x\mapsto m]\big),\\
\sem{P\mid Q}\rho&=\sem{P}\rho\mid\sem{Q}\rho,&
\sem{\pv{new}\ a;P}\rho&=\anew.\big(m\mapsto\sem{P}\rho[a\mapsto m]\big),\\
\sem{!P}\rho&=\,!\sem{P}\rho,&
\sem{\pv{let}\ x=D\ \pv{in}\ P\ \pv{else}\ Q}\rho&=
\begin{cases}\sem{P}\rho[x\mapsto m]&\sem{D}\rho=m\in\Msg,\\ \sem{Q}\rho&\sem{D}\rho=\bot,\end{cases}\\
&&\sem{\pv{if}\ M\ \pv{then}\ P\ \pv{else}\ Q}\rho&=
\begin{cases}\sem{P}\rho&\rho M=\true,\\ \sem{Q}\rho&\rho M\ne\true.\end{cases}
\end{align*}
The translation of $P_0$ is $\sem{P_0}\rho_0$. A translated process mentions no $\nu_k$ except through $\anew$.
\end{definition}

\begin{remark}
In the development a process is a program in the continuation monad with answer type $\Proc$, run by $\mathrm{run}(b)=b(\lambda\_.\,\mathbf{0})$, with primitives $\mathrm{send}(c,m)=\lambda\kappa.\,\aout(c,m).\kappa()$, $\mathrm{recv}(c)=\lambda\kappa.\,\ain(c).\kappa$, $\mathrm{fresh}=\lambda\kappa.\,\anew.\kappa$, $\mathrm{emit}(m)=\lambda\kappa.\,\aev(m).\kappa()$, $\mathrm{fork}=\lambda\kappa.\,\kappa(\mathrm{tt})\mid\kappa(\mathrm{ff})$, $\mathrm{repl}=\lambda\kappa.\,!\kappa()$ and $\mathrm{stop}=\lambda\_.\,\mathbf{0}$; $\pv{let}$ and $\pv{if}$ are case analysis in the metalanguage. The clauses above are the unfolding of that presentation.
\end{remark}

\subsection{The attacker and the security properties}

\begin{definition}[Enemy]
The processes $\Enemy(K)\in\Proc$, for $K\subseteq\Msg$, are the pointwise least such that
\begin{enumerate}
\item $\varepsilon\in\Enemy(K)$;
\item if $K\vdash c$ and $e\in\Enemy(K\cup\{m\})$ then $\ain(c,m)\,e\in\Enemy(K)$, for every $m\in\Msg$;
\item if $K\vdash c$, $K\vdash m$ and $e\in\Enemy(K)$ then $\aout(c,m)\,e\in\Enemy(K)$;
\item if $K\vdash c$, $K\vdash m$ and $e\in\Enemy(K)$ then $\acomm(c,m)\,e\in\Enemy(K)$, an output to itself;
\item if $k\notin\nonces(e)$ and $e\in\Enemy(K\cup\{\nu_k\})$ then $\anew(k)\,e\in\Enemy(K)$.
\end{enumerate}
\end{definition}

\begin{definition}[Attacks, knowledge]
For $P\in\Proc$,
\[
\Attacks(P)=\{t\in P\mid\Enemy(\emptyset)\ :\ \Closed(t)\}.
\]
For $K\subseteq\Msg$, $\known(K,\cdot):\Act^*\to\mathcal{P}(\Msg)$ is defined by
\begin{align*}
\known(K,\varepsilon)&=K,\\
\known(K,a\,t)&=\known\big(\langle K\cup\{m\mid a=\acomm(c,m),\ c\in K\}\rangle,\,t\big).
\end{align*}
\end{definition}

\begin{proposition}[Prefix-closure]\label{prop:prefix}
$\mathbf{0}$ and $\Enemy(K)$ are prefix-closed. If $P$, every $F(m)$ and every $P(i)$ are prefix-closed, so are $a\cdot P$, $\ain(c).F$, $\anew.F$ and $\mrg_I(P)$. If $P$ is prefix-closed, so is $\Attacks(P)$.
\end{proposition}

\begin{definition}[Facts, secrecy, correspondence]
A \emph{fact} is $\event(e)$ or $\knows(m)$ with $e,m\in\Msg$; $\Fact$ is the set of facts. A fact \emph{holds} after $t\in\Act^*$, written $t\models f$, if
\[
t\models\event(e)\iff\aev(e)\in t,\qquad
t\models\knows(m)\iff m\in\known(\langle\emptyset\rangle,t).
\]
For $F,H\in\Fact$ and $s\in\Msg$,
\[
\Corr(P,F,H)\iff\forall t\in\Attacks(P).\ t\models F\Rightarrow t\models H,\qquad
\Secret(P,s)\iff\forall t\in\Attacks(P).\ t\not\models\knows(s).
\]
By Proposition~\ref{prop:prefix}, for prefix-closed $P$ these are properties of every prefix of every attack: $\Corr(P,F,H)$ says that $H$ holds by the time $F$ does.
\end{definition}

\begin{definition}[Queries]\label{def:queries}
A \emph{query fact} is $\event(M)$ or $\attacker(M)$ with $M$ a term over constructors, free names and a finite set $X$ of variables; it is an \emph{event fact} or an \emph{attacker fact} accordingly. For a valuation $\sigma:X\to\Msg$ write $\sigma M\in\Msg$ for the evaluation of $M$ with free names as constants, and let the \emph{instance} of a query fact under $\sigma$ be $\sigma\,\event(M)=\event(\sigma M)$ and $\sigma\,\attacker(M)=\knows(\sigma M)$.
The query $\pv{query}\ \attacker(M)$, for $M$ without variables, denotes $\Secret(\sem{P_0}\rho_0,\rho_0M)$. The query $\pv{query}\ F\Longrightarrow H$, with $F$ and $H$ query facts over $X$ such that every variable of $H$ occurs in $F$, denotes $\forall\sigma:X\to\Msg.\ \Corr(\sem{P_0}\rho_0,\sigma F,\sigma H)$.
\end{definition}

\begin{definition}[Bounded traces]
For $S\subseteq\Msg$, \emph{$S$-boundedness} of a trace is defined by recursion: $\varepsilon$ is $S$-bounded; $\ain(c,m)\,u$ is $S$-bounded if $u$ is whenever $m\in S$; $\aout(c,m)\,u$ and $\acomm(c,m)\,u$ are $S$-bounded if $m\in S$ and $u$ is; and $a\,u$ is $S$-bounded if $u$ is, for every other action $a$.
\end{definition}

\begin{lemma}\label{lem:enemy}
Let $S\subseteq\Msg$ be $\Ops$-closed with $\nu_k\in S$ for all $k$ and $K\subseteq S$. Every $e\in\Enemy(K)$ is $S$-bounded.
\end{lemma}

\begin{lemma}\label{lem:network}
Let $\comms(s,t)$ with $\Closed(t)$, and let $s{\restriction}i$ be $S$-bounded for every $i$. Then $m\in S$ for every $\acomm(c,m)$ in $t$.
\end{lemma}

\begin{lemma}\label{lem:known}
Let $S\subseteq\Msg$ be $\Ops$-closed, $K\subseteq S$, and $m\in S$ for every $\acomm(c,m)$ in $t$. Then $\known(K,t)\subseteq S$.
\end{lemma}

\begin{theorem}[Closed-set method]\label{thm:closed}
Let $P\in\Proc$ and let $S\subseteq\Msg$ be $\Ops$-closed with $\nu_k\in S$ for all $k$ and such that every trace of $P$ is $S$-bounded. Then $\known(\langle\emptyset\rangle,t)\subseteq S$ for every $t\in\Attacks(P)$; in particular $s\notin S$ implies $\Secret(P,s)$.
\end{theorem}

\section{Relation to the operational semantics of ProVerif}

\subsection{The reduction semantics}

This section restates the semantics of Blanchet's survey (\S2.4, \S2.5.1, \S3.1--3.3) for the syntax of Definition~\ref{def:syntax}. Names are atoms; substitutions map variables to terms and never substitute $\fail$ or a may-fail variable for an ordinary variable. An equational theory $E$ is a finite set of equations between terms without names; $=_E$ is the closure of $E$ under reflexivity, symmetry, transitivity, application of function symbols and substitution. Each constructor $f$ carries the rule $f(\vec x)\to f(\vec x)$, each destructor $g$ a finite ordered list of rules $g(\vec U)\to U$ between may-fail terms without names, and every $\mathrm{def}(h)$ ends with $h(\vec u)\to\fail$ for may-fail variables $\vec u$.

Evaluation $D\ev U$ of closed expressions to closed may-fail terms is the least relation with $M\ev M$, $\fail\ev\fail$, and $h(D_1,\dots,D_n)\ev\sigma U'_j$ whenever $D_i\ev U_i$ for all $i$, $\mathrm{def}(h)=(h(U'_{i,1},\dots,U'_{i,n})\to U'_i)_{i\le k}$, $\sigma U'_{j,i}=_E U_i$ for all $i$, and for all $i<j$ and all $\sigma'$ there is some $l$ with $\sigma'U'_{i,l}\ne_E U_l$.

A \emph{configuration} is $(\Npub,\Npriv),\mathcal{P}$ with $\Npub,\Npriv$ finite disjoint sets of names and $\mathcal{P}$ a finite multiset of closed processes with $\fn(\mathcal{P})\subseteq\Npub\cup\Npriv$. Writing $\mathcal{E}$ for the pair of name sets, reduction is the least relation closed under
\begin{align*}
\mathcal{E},\mathcal{P}\cup\{\mathbf{0}\}&\to\mathcal{E},\mathcal{P} \tag{Red Nil}\\
\mathcal{E},\mathcal{P}\cup\{P\mid Q\}&\to\mathcal{E},\mathcal{P}\cup\{P,Q\} \tag{Red Par}\\
\mathcal{E},\mathcal{P}\cup\{!P\}&\to\mathcal{E},\mathcal{P}\cup\{P,!P\} \tag{Red Repl}\\
(\Npub,\Npriv),\mathcal{P}\cup\{\pv{new}\ a;P\}&\to(\Npub,\Npriv\cup\{a'\}),\mathcal{P}\cup\{P\{a'/a\}\}\notag\\
&\qquad\text{if } a'\notin\Npub\cup\Npriv \tag{Red Res}\\
\mathcal{E},\mathcal{P}\cup\{\pv{out}(N,M);Q,\ \pv{in}(N',x);P\}&\to\mathcal{E},\mathcal{P}\cup\{Q,\ P\{M/x\}\}\quad\text{if } N=_E N' \tag{Red I/O$'$}\\
\mathcal{E},\mathcal{P}\cup\{\pv{let}\ x=D\ \pv{in}\ P\ \pv{else}\ Q\}&\to\mathcal{E},\mathcal{P}\cup\{P\{M/x\}\}\quad\text{if } D\ev M \tag{Red Eval 1}\\
\mathcal{E},\mathcal{P}\cup\{\pv{let}\ x=D\ \pv{in}\ P\ \pv{else}\ Q\}&\to\mathcal{E},\mathcal{P}\cup\{Q\}\quad\text{if } D\ev\fail \tag{Red Eval 2}\\
\mathcal{E},\mathcal{P}\cup\{\pv{if}\ M\ \pv{then}\ P\ \pv{else}\ Q\}&\to\mathcal{E},\mathcal{P}\cup\{P\}\quad\text{if } M=_E\true \tag{Red Cond 1$'$}\\
\mathcal{E},\mathcal{P}\cup\{\pv{if}\ M\ \pv{then}\ P\ \pv{else}\ Q\}&\to\mathcal{E},\mathcal{P}\cup\{Q\}\quad\text{if } M\ne_E\true \tag{Red Cond 2$'$}\\
\mathcal{E},\mathcal{P}\cup\{\pv{event}(M);P\}&\to\mathcal{E},\mathcal{P}\cup\{P\} \tag{Red Event}
\end{align*}
A \emph{trace} is a finite reduction sequence $\mathcal{E}_0,\mathcal{P}_0\to\dots\to\mathcal{E}_n,\mathcal{P}_n$, written $\mathcal{E}_0,\mathcal{P}_0\to^*\mathcal{E}_n,\mathcal{P}_n$.

\begin{definition}[Survey, Definitions 3.1--3.5]\label{def:survey}
Let $P_0$ be a closed process.
\begin{enumerate}
\item A closed process $Q$ is an \emph{$\Npub$-adversary} if $\fn(Q)\subseteq\Npub$, every function symbol of $Q$ is public, and $Q$ contains no $\pv{event}$.
\item A trace \emph{outputs $M$ publicly} if it contains a step \red{I/O$'$} whose output is $\pv{out}(N,M);Q$ with $N\in\Npub$. It \emph{executes $\pv{event}(M)$} if it contains a step \red{Event} on $\pv{event}(M);P$.
\item $P_0$ \emph{preserves the secrecy of} the closed term $M$ \emph{from} $\Npub$ if, for some $\Npriv$ disjoint from $\Npub$ with $\fn(M)\cup\fn(P_0)\subseteq\Npub\cup\Npriv$, no trace $(\Npub,\Npriv),\{P_0,Q\}\to^*\mathcal{E}',\mathcal{P}'$ with $Q$ an $\Npub$-adversary outputs $M$ publicly.
\item $P_0$ \emph{satisfies} $\event(M)\rightsquigarrow\bigvee_{j=1}^m\bigwedge_{k=1}^{l_j}\event(M_{jk})$ \emph{against $\Npub$-adversaries} if, for some $\Npriv$ as above, for every $\Npub$-adversary $Q$, every trace $\mathrm{Tr}$ from $(\Npub,\Npriv),\{P_0,Q\}$ and every substitution $\sigma$: if $\mathrm{Tr}$ executes $\pv{event}(\sigma M)$ then there are $\sigma'$ and $j$ with $\sigma'M=\sigma M$ such that $\mathrm{Tr}$ executes $\pv{event}(\sigma'M_{jk})$ for every $k\le l_j$.
\end{enumerate}
\end{definition}

Fix $\Npriv$ as in Definition~\ref{def:survey}, and for an $\Npub$-adversary $Q$ write $\CQ$ for the configuration $(\Npub,\Npriv),\{P_0,Q\}$.

\begin{definition}[Satisfaction]\label{def:satisfies}
A trace \emph{satisfies} $\attacker(M)$ if it outputs $M$ publicly, and \emph{satisfies} $\event(M)$ if it executes $\pv{event}(M)$.
\end{definition}

\begin{definition}[Correspondences between facts]\label{def:corr-facts}
Let $F$ and $H_{jk}$ ($j\le m$, $k\le l_j$) be query facts over $X$. $P_0$ \emph{satisfies} $F\rightsquigarrow\bigvee_{j=1}^m\bigwedge_{k=1}^{l_j}H_{jk}$ \emph{against $\Npub$-adversaries} if for every $\Npub$-adversary $Q$, every trace $\mathrm{Tr}$ from $\CQ$ and every substitution $\sigma$: if $\mathrm{Tr}$ satisfies $\sigma F$ then there are $\sigma'$ and $j$ with $\sigma'F=\sigma F$ such that $\mathrm{Tr}$ satisfies $\sigma'H_{jk}$ for every $k\le l_j$. For event facts this is Definition~\ref{def:survey}(4).
\end{definition}

\subsection{The correspondence}

Throughout, fix a model (Definition~\ref{def:model}) satisfying
\begin{enumerate}
\item[(R1)] $\Fr$ is the set of symbols at the root of no equation and $(E,\Fr)$ satisfies (H);
\item[(R2)] every rule of every destructor has patterns over $\Fr$, and the destructor is the one of Definition~\ref{def:rules};
\item[(R3)] $\Pub$ consists of the public constructors and the names in $\Npub$, and $\Dt$ of the public destructors;
\item[(R4)] $P_0$ is a closed process of Definition~\ref{def:syntax}.
\end{enumerate}
Names of the reduction semantics are the free names together with the created names; terms are interpreted in $\Msg$ through an injection $\iota$ from created names to $\N$, extended by $\iota a=a$ for free names, $\iota a'=\nu_{\iota a'}$ for created names, $\iota f(\vec M)=f(\iota\vec M)$ and $\iota\,\fail=\bot$; $\iota M$ depends only on the class of $M$ modulo $E$. Let $\rho_\iota$ be the environment with $\rho_\iota(a)=\iota a$ for every free or created name $a$.

\begin{lemma}[Evaluation]\label{lem:eval}
For every closed expression $D$ over free and created names, some $U$ satisfies $D\ev U$, and $D\ev U$ implies $\iota U=\sem{D}\rho_\iota$.
\end{lemma}

\begin{definition}[Labelled semantics]
A \emph{labelled configuration} is $(U,\mathcal{P})$ with $U\subseteq\N$ finite and $\mathcal{P}$ a finite multiset of closed processes whose names are free names and symbols $\nu_k$. The $\tau$-steps $(U,\mathcal{P})\xrightarrow{\tau}(U,\mathcal{P}')$ are the steps \red{Nil}, \red{Par}, \red{Repl}, \red{Eval 1}, \red{Eval 2}, \red{Cond 1$'$} and \red{Cond 2$'$}, with $=_E$ read on $\Msg$. The visible steps are
\begin{align*}
(U,\mathcal{P}\cup\{\pv{out}(N,M);P\})&\xrightarrow{\aout([N],[M])}(U,\mathcal{P}\cup\{P\}),\\
(U,\mathcal{P}\cup\{\pv{in}(N,x);P\})&\xrightarrow{\ain([N],[M])}(U,\mathcal{P}\cup\{P\{M/x\}\})\quad\text{for every closed term }M,\\
(U,\mathcal{P}\cup\{\pv{out}(N,M);Q,\ \pv{in}(N',x);P\})&\xrightarrow{\acomm([N],[M])}(U,\mathcal{P}\cup\{Q,\ P\{M/x\}\})\quad\text{if }N=_E N',\\
(U,\mathcal{P}\cup\{\pv{new}\ a;P\})&\xrightarrow{\anew(k)}(U\cup\{k\},\mathcal{P}\cup\{P\{\nu_k/a\}\})\quad\text{for every }k\notin U,\\
(U,\mathcal{P}\cup\{\pv{event}(M);P\})&\xrightarrow{\aev([M])}(U,\mathcal{P}\cup\{P\}).
\end{align*}
The third step is \red{I/O$'$} between two members of $\mathcal{P}$. It is visible rather than a $\tau$-step because a trace records a completed communication, which is what $\known$ reads the attacker's knowledge off. With $\xRightarrow{\alpha}\;=\;\xrightarrow{\tau}^{*}\xrightarrow{\alpha}\xrightarrow{\tau}^{*}$, the set of \emph{weak traces} is
\[
\Tr(U,\mathcal{P})=\{\alpha_1\cdots\alpha_n\in\Act^*\mid (U,\mathcal{P})\xRightarrow{\alpha_1}\cdots\xRightarrow{\alpha_n}\}.
\]
\end{definition}

\begin{theorem}[Adequacy]\label{thm:T}
For every closed process $P$ whose names are free names, $\sem{P}\rho_0=\Tr(\emptyset,\{P\})$. For all finite multisets $\mathcal{P},\mathcal{Q}$ of such processes, $\Tr(\emptyset,\mathcal{P}\uplus\mathcal{Q})=\Tr(\emptyset,\mathcal{P})\mid\Tr(\emptyset,\mathcal{Q})$.
\end{theorem}

\begin{definition}[Observation]
For a trace $\mathrm{Tr}$ from $\CQ$, $\obs(\mathrm{Tr})\in\Act^*$ is the word that records $\anew(\iota a')$ for each step \red{Res} creating $a'$ and $\aev(\iota M)$ for each step \red{Event} on $\pv{event}(M)$, in the order of the steps, and nothing for the other steps. For $t\in\Act^*$, $t{\restriction}_{\anew,\aev}$ is the subword of $\anew$ and $\aev$ actions.
\end{definition}

Let $\Azero=\Attacks(\sem{P_0}\rho_0)$.

\begin{theorem}[Attacks]\label{thm:A}
The following hold.
\begin{enumerate}
\item[(S)] For every $\Npub$-adversary $Q$ and every trace $\mathrm{Tr}$ from $\CQ$ there is $t\in\Azero$ with $t{\restriction}_{\anew,\aev}=\obs(\mathrm{Tr})$ such that $\iota M\in\known(\langle\emptyset\rangle,t)$ whenever $\mathrm{Tr}$ outputs $M$ publicly.
\item[(C)] For every $t\in\Azero$ there are an $\Npub$-adversary $Q$, a trace $\mathrm{Tr}$ from $\CQ$ and an injection $\iota$ with $\obs(\mathrm{Tr})=t{\restriction}_{\anew,\aev}$; if moreover $m\in\known(\langle\emptyset\rangle,t)$ and $\Npub\ne\emptyset$, $\mathrm{Tr}$ extends to a trace with the same observation whose last step outputs publicly a term $M$ with $\iota M=m$.
\end{enumerate}
\end{theorem}

\begin{theorem}[Secrecy]\label{thm:secrecy}
Let $\Npub\ne\emptyset$, let $M$ be a closed term with $\fn(M)\subseteq\Npub\cup\Npriv$ and let $s=\rho_0M$. Then $\Secret(\sem{P_0}\rho_0,s)$ holds if and only if $P_0$ preserves the secrecy of $M$ from $\Npub$.
\end{theorem}

\begin{remark}
The hypothesis $\Npub\ne\emptyset$ cannot be dropped: when $\Npub=\emptyset$ no trace outputs anything publicly, so $P_0$ preserves the secrecy of every $M$, whereas $\Secret(\sem{P_0}\rho_0,\rho_0M)$ fails whenever $\emptyset\vdash\rho_0M$, for instance for a public constant $M$.
\end{remark}

\begin{theorem}[Correspondences]\label{thm:corr}
Let $F$ and $H$ be query facts over $X$ such that $H$ is an event fact and $\sigma F=\sigma'F$ implies $\sigma H=\sigma'H$ for all valuations $\sigma,\sigma'$, and let $\Npub\ne\emptyset$ if $F$ is an attacker fact. Then $P_0$ satisfies $F\rightsquigarrow H$ against $\Npub$-adversaries if and only if $\Corr(\sem{P_0}\rho_0,\sigma F,\sigma H)$ holds for every $\sigma:X\to\Msg$.
\end{theorem}

\begin{remark}
The hypotheses on $H$ and on $\Npub$ cannot be dropped. When $\Npub=\emptyset$ and $F=\attacker(M)$, $P_0$ satisfies the correspondence vacuously, whereas $\Corr$ fails as soon as $\emptyset\vdash\sigma M$ for some valuation $\sigma$. For $H=\attacker(M)$, Definition~\ref{def:satisfies} requires $\sigma M$ itself to be output publicly, whereas $\knows(\sigma M)$ holds as soon as $\sigma M$ is derivable: with the symmetric encryption of the survey's \S2.1, $\Npub=\{c\}$ and $P_0=\pv{out}(c,\pv{senc}(s,k));\pv{out}(c,k);\pv{event}(e)$, $\Corr$ holds for $\event(e)\rightsquigarrow\attacker(s)$ while the correspondence of Definition~\ref{def:corr-facts} fails. The reading by derivability is ProVerif's.
\end{remark}

\end{document}
